\section{Birefringence of a calcite crystal}

    \subsection{Theory of birefringence}
    Birefringence, also known as double refraction, is a phenomenon that occurs in certain anisotropic materials, such as calcite crystals.
    When light enters a birefringent material, it splits into two rays: the ordinary ray and the extraordinary ray. These rays experience different
    refractive indices due to the anisotropic nature of the material, resulting in different propagation speeds and directions. The ordinary ray follows
    Snell's law, while the extraordinary ray does not, leading to a separation of the two rays within the crystal.

    \subsection{Experimental setup for birefringence}
    In this experiment, a calcite crystal was placed in the path of an unpolarized light beam. As the light passed through the calcite crystal, it split into
    the ordinary and extraordinary rays. A analyzer was placed behind the crystal to observe the polarization states of the two rays. The rays were projected
    onto a screen. By rotating the polarizer and observing the changes in intensity and direction of the rays, we can analyze the birefringent properties of
    the calcite crystal.

    \subsection{Results of birefringence}
    